\scanpage{117}
\mspage{348}

Here it is clear that the only values modulo $3$ \unsure{represented} by the quadratic form $3c^2+3cd+d^2\equiv d^2\ (3)$ are $0$ and $1$, whence
\[
(102)\quad p\equiv-1\ (3)\quad(\text{i.e. }\gamma\equiv-1\ (3))\quad\text{if }\xi=\lambda l_0^{-1}\rho(\lambda)^{-1}
\]
\[
(\Longleftrightarrow g=\lambda_0l_0^{-1}\rho(\lambda_0)^{-1})
\]

This study thus does show that we are in the case $\xi=\lambda\rho(\lambda)^{-1}$, resp. $\xi=\lambda l_0^{-1}\rho(\lambda)^{-1}$, according as $p=\chi(u)\equiv1$ or $\equiv-1\ (3)$, i.e. according as $\gamma\equiv0$ or $\gamma\equiv-1\ (3)$.

\uline{Remarks.} 1°) \unsure{Placing} ourselves in $\mathrm{Sl}(2,\widehat{\mathbb Z})'$ \deleted{or else} \inserted{(or rather, in its} \deleted{of the $\unsure{\Gamma|F}$} subgroup $\Pi(2)'=\operatorname{Ker}(\mathrm{Sl}(2,\widehat{\mathbb Z})'\unsure{-|\longrightarrow}\mathrm{Sl}(2,\mathbb Z/2\mathbb Z))$ in place of $\widehat{\Pi}_{0,3}$, to \deleted{\unsure{solve the equation}} \inserted{determine the} solutions of ``the equation in $\xi$'' (**) (cf. (68)), when $\xi$ is put in \unsure{parametric} form (94), we see that as a condition on $\lambda=\begin{pmatrix}a&b\\c&d\end{pmatrix}$, this equation is relatively innocuous and very simple --- it simply says that $pD=1$, where $D=c^2+d^2+cd$ in the first case, $D=\deleted{-(3c^2+3cd+d^2)}\inserted{-(c^2+d^2+cd)}$ in the second. If we do not fix $p\in\widehat{\mathbb Z}^{*}$ beforehand, this condition is simply that the element

\scanpage{118}

$D$ of $\widehat{\mathbb Z}$ be a unit, i.e.
\[
(103)\quad D\in\widehat{\mathbb Z}^{*}
\]
--- and we then determine $p$ uniquely in terms of $D$ \inserted{, hence of $c$ and $d$ and the sign $\delta$ by $D=\delta(c^2+d^2+cd)$, \deleted{and of $\delta\in\{\pm1\}$}} ($a$ and $b$ are of no use) --- hence also in terms of $\lambda$\inserted{$\in\widehat{\Pi}_{0,3}\cdot\unsure{(1,\tau)|\{1,\tau\}}$}, itself uniquely determined by $\xi$.
\margin{\uline{N.B.} We \unsure{will have} $p$}

\uline{Thus knowledge of} $\xi\in\widehat{\Pi}_{0,3}$ \uline{suffices} \uline{to determine} \inserted{\uline{the multiplier}} $p=2\gamma+1$\inserted{$\in\widehat{\mathbb Z}^{*}$}, hence also $\gamma\in\widehat{\mathbb Z}$, and $h_0=l_0^{-\gamma}\xi$, hence also $g_0$ (by formula (66), $g_0^{-1}\sigma_\infty(g_0)=h_0$), and consequently $g$ by formula (62), $g=\sigma_\infty(g_0)l_1^\gamma\rho(g_0)^{-1}$, hence $u$ by (50). It is \unsure{interesting} to make sure that \unsure{$D$|it} \deleted{\illegible{1}} does not change when we make the change (69)
\[
\xi\longmapsto\xi'=l_0^{-\alpha}\xi l_1^\alpha,
\]
in matrix terms, we have indeed
\[
(104)\quad l_0^{-\alpha}\xi l_1^\alpha=
\begin{pmatrix}1&2\alpha\\0&1\end{pmatrix}
\begin{pmatrix}A&B\\C&D\end{pmatrix}
\begin{pmatrix}1&0\\2\alpha&1\end{pmatrix}
=\begin{pmatrix}A+2\alpha(B+C)+4\alpha^2D&B+2\alpha D\\C+2\alpha D&D\end{pmatrix}.
\]

2/ But we should take care that in \unsure{full} \uline{rigor} mere knowledge of \uline{the image of $\xi$ in} $\mathrm{Sl}(2,\widehat{\mathbb Z})'$ allows us to determine the multiplier $p$ \uline{only up to sign} --- on the other hand knowledge of the image of $\lambda$ in $\mathrm{Gl}(2,\widehat{\mathbb Z})'$ --- hence of a matrix $\begin{pmatrix}a&b\\c&d\end{pmatrix}\in\unsure{\mathrm{Sl}|\mathrm{Gl}}(2,\widehat{\mathbb Z})$ defined up to \uline{sign} --- allows us to determine

\scanpage{119}
\mspage{349}

$p$ as the inverse of $D$, given \deleted{depending on the case} by $D=\delta(c^2+cd+d^2)$ \deleted{\unsure{or} $D=3(c^2+cd)+d^2$}. \inserted{where $\delta\in\{\pm1\}$ is determined by the condition $D\equiv\delta\ (3)$.}

3/ We know by arithmetic means (since $\Pi_{\mathbb Q}\xrightarrow{\unsure{\chi}}\widehat{\mathbb Z}^{*}$ is surjective) that we must find any $p\in\widehat{\mathbb Z}^{*}$ in the form $1/D$ (or again, in the form $D$), where $D=c^2+cd+d^2$ if $p\equiv1\ (3)$, and $D=-3(c^2+cd)\unsure{-|+}d^2$ if $p\equiv-1\ (3)$, with $c,d\in\widehat{\mathbb Z}$, and even $c\equiv0\ (2)$, $d\equiv1\ (2)$. \inserted{If we want the direct verification,} (We are reduced to verifying it separately for the various $\mathbb Z_\ell$, i.e. to proving the

\uline{Corollary.} \unsure{1°)|°)} If $\ell$ \deleted{\illegible{1}} prime $\ne3$, then $\forall\unsure{\beta|p}\in\mathbb Z_\ell^{*}$ can be put equally well in the forms $c^2+cd+d^2$, or $-(3(c^2+cd)+d^2)$ (with $c,d\in\unsure{\widehat{\mathbb Z}_\ell|\mathbb Z_\ell}$). If $\ell=2$, we can even suppose $c\equiv0$, $d\equiv1\ (2)$. For $p\in\mathbb Z_\ell^{*}$, $\ell=3$, it can be put in the form $c^2+cd+d^2$ (resp. $-3(c^2+cd)\unsure{-|+}d^2$) iff $p\equiv1\ (3)$ (resp. $p\equiv-1\ (3)$), cases which are mutually exclusive.
\margin{we don't care \unsure{about the|about} $-(3(c^2+cd)+d^2)$ !}

4/ \unsure{Returning} to the ``innocuous'' reduction in $\mathrm{Sl}(2,\widehat{\mathbb Z})'$ of ``the equation in $\lambda$'' as meaning simply that $c^2+d^2+cd\in\widehat{\mathbb Z}^{*}$ \deleted{resp. $3(c^2+cd)+d^2\in\widehat{\mathbb Z}^{*}$}, we might hope

\scanpage{120}

to have a method for constructing even in $\widehat{\Pi}_{0,3}$ big batches of solutions of the equation in question. But I have absolutely not managed it --- not even to construct a single one explicitly, apart from \unsure{those} which correspond to the case $u=1$, and $u=\tau_\infty$.

A question: Let $\lambda_0\in\mathrm{Sl}(2,\widehat{\mathbb Z})^{\unsure{\prime|*}}$, $\lambda_0\equiv1\ (2)$, be such that $\xi_0=\lambda_0\rho(\lambda_0)^{-1}$ (say, or $\xi=\lambda(-l_0^{-1})\rho(\lambda_0)^{-1}$) \deleted{be such that} \unsure{satisfies} in $\mathrm{Sl}(2,\widehat{\mathbb Z})'$ an ``equation in $\xi$'' for a suitable $\gamma=\frac{p-1}{2}$ ($p\in\widehat{\mathbb Z}^{*}$) --- i.e. $c^2+d^2+cd\in\widehat{\mathbb Z}^{*}$ (resp. $3(c^2+cd)+d^2\in\widehat{\mathbb Z}^{*}$). Does there then exist a $\unsure{\lambda|\xi}\in\widehat{\Pi}_{0,3}$ above $\lambda_0$, such that $\xi=\lambda\rho(\lambda)^{-1}$ (resp.\ldots) \unsure{satisfies} an equation in $\xi$ as well (\unsure{necessarily} with the same $\gamma$) ?

\scanpage{121}
\mspage{350}

\uline{43} \deleted{Return to the notations} -- Digression on \inserted{the 1-}cocycles of cyclic groups.

\unsure{Let|Let} $\Pi$ be a group, $\rho$ an \unsure{automorphism} of $\Pi$ such that
\[
\rho^n=\mathrm{id} \tag{1}
\]
where $n\in\mathbb N^*$. We thus have an operation of the cyclic group $\mathbb Z/n\mathbb Z$ on $\Pi$, whence a semidirect product
\[
\widetilde\Pi=\Pi.\,\mathbb Z/n\mathbb Z, \tag{\unsure{1|2}}
\]
and an extension structure

(3) \diagram{\begin{tikzcd}[ampersand replacement=\&]
1 \arrow[r] \& \Pi \arrow[r] \& \widetilde\Pi \arrow[r] \& \mathbb Z/n\mathbb Z \arrow[r] \arrow[l,bend left=30,dotted] \& 1.
\end{tikzcd}}

We \deleted{identify}\inserted{denote by} $\dot\rho$ \deleted{\unsure{with}}\inserted{the} generator of $\mathbb Z/n\mathbb Z$, and \unsure{we|to} \inserted{\unsure{we identify $\rho$ with}}\deleted{\illegible{1}} the image of the latter in $\widetilde\Pi$, by the canonical splitting \deleted{\illegible{1}}
\[
\mathbb Z/n\mathbb Z\longrightarrow\widetilde\Pi\qquad 1\bmod n\longmapsto\rho. \tag{4}
\]
We thus have
\[
\rho(g)=\rho g\rho^{-1}\qquad(g\in\Pi). \tag{5}
\]
The splittings of the extension (3) correspond bijectively to the elements $\rho'$ of $\widetilde\Pi$ over $\dot\rho$, such that
\[
\rho'{}^n=1. \tag{6}
\]
To say that $\rho'$ is over $\unsure{\rho|\dot\rho}$ means that we have
\[
\rho'=\rho'_g\overset{\unsure{\mathrm{dfn}}}{=}g\rho\qquad\text{with }g\in\Pi, \tag{7}
\]
and in terms of such a $g\in\Pi$ equation (6) \unsure{Is|is} rewritten as follows
\[
\boxed{g\rho(g)\rho^2(g)\cdots\rho^{n-1}(g)=1} \tag{8}
\]

\scanpage{122}

\deleted{Let us set, for $g\in\Pi$}

\deleted{We will have, for $\alpha\in\Pi$,}
\deleted{$\unsure{\alpha\rho\alpha^{-1}=}\rho'_g=g\rho$}

\deleted{so that \illegible{2}}

\deleted{We will have, for \illegible{1} $g$ as above.}

The group $\Pi$ operates by conjugation on the \unsure{set} of splittings of the extension (3), or what amounts to the same thing, on the \unsure{set} of $\rho'$ above $\unsure{\rho|\dot\rho}$, satisfying (6), \deleted{We} We will moreover have, for $\alpha\in\Pi$
\[
\begin{aligned}
\alpha\rho'_g\alpha^{-1}
&=\deleted{\unsure{(\alpha g\alpha^{-1})(\alpha\rho\alpha^{-1})=\alpha}}\inserted{\alpha g\rho\alpha^{-1}=\alpha g\rho\alpha^{-1}\rho^{-1}\rho}\\
&=(\alpha g\rho(\alpha)^{-1})\rho
=\unsure{\rho'_{\alpha g\rho(\alpha)^{-1}}}
\end{aligned}\tag{9}
\]
i.e. the operation by conjugation of $\Pi$ on the \unsure{set} of splittings of (3), is interpreted as the operation $\Theta=\Theta_\rho$ of $\Pi$ on the \unsure{set} of solutions of (8) \inserted{given, for $\alpha\in\Pi$, by}
\[
\unsure{\Theta_\rho(\alpha)}:g\longmapsto\alpha g\rho(\alpha)^{-1} \tag{10}
\]
(We will say that $\alpha g\rho(\alpha)^{-1}$ is the $\rho$-conjugate of $g$ by the element $\alpha\in\Pi$). The $\Pi$-conjugacy classes \unsure{of|of the} splittings of (3) thus correspond to the $\Pi$-$\rho$-conjugacy classes of solutions of (8).

Among these classes, the \inserted{``trivial''} class \deleted{\unsure{of the extension}}\inserted{\unsure{of the splitting}} canonical for (3), i.e. the class of $g=1$, is given parametrically \deleted{by} as the \unsure{set} of
\[
g=\alpha\rho(\alpha)^{-1}\qquad(\alpha\in\Pi). \tag{11}
\]
Let $g,g'$ be two \unsure{congruent|conjugate} solutions of

\scanpage{123}
\mspage{351}

the equation (8), so that we have
\[
g'=\alpha g\rho(\alpha)^{-1}\qquad\text{i.e. }g'=\Theta(\alpha).g \tag{12}
\]
for a suitable $\alpha\in\Pi$. When do we have uniqueness \unsure{of the|of} $\alpha$ -- more generally, what can we say, for fixed $g,g'$, of the \unsure{set} of $\alpha\in\Pi$ \deleted{which} such that we have (12) (which we can call the $\rho$-\uline{transporter} of $g$ into $g'$, in the group $\Pi$). \deleted{As}\inserted{It is} obviously a right coset \unsure{modulo|module} the subgroup $\Pi(\rho,g)$ of $\Pi$ which \deleted{stabilizes} stabilizes $g$
\[
\begin{aligned}
\Pi(\rho,g)&=\{\beta\in\Pi\mid\Theta_\rho(\beta).g=g\}\\
&=\{\beta\in\Pi\mid\beta g\rho(\beta)^{-1}=g\}\\
&=\{\beta\in\Pi\mid\unsure{\{\mathrm{int}(\beta)\}}(g\rho)=g\rho\}\\
&=\Pi^{\deleted{g}\rho'_g}=\{\beta\in\Pi\mid\mathrm{int}(\underset{\inserted{\rho'_g}}{\uline{g\rho}})\beta=\beta\}\\
&=\text{Centralizer of }\deleted{g}\rho'_g\ \uline{\text{in }\Pi}.
\end{aligned}\tag{13}
\]

In other words, if $\alpha$ satisfies (12), the \unsure{set} of $\alpha'$ which satisfy the same relation is none other than the \unsure{set} $\alpha\Pi(\rho,g)$ :
\[
\begin{aligned}
\alpha'g\rho(\alpha')^{-1}=\alpha g\rho(\alpha)^{-1}
&\Longleftrightarrow\alpha'\in\alpha\Pi(\rho,g)=\alpha\Pi^{\rho'_g}\\
&\text{i.e. }\alpha^{-1}\alpha'\in\Pi^{\rho'_g}
\end{aligned}\tag{14}
\]
Thus we have uniqueness of the $\alpha$ satisfying (12) \deleted{\unsure{iff}} \unsure{iff}
\[
\Pi^{\rho'_g}=1, \tag{15}
\]
a condition which depends only on the conjugacy class of the splitting \deleted{\unsure{under consideration}} of (3) defined by $g$ i.e. by $\rho'_g$.

\scanpage{124}

In fact, not only the group $\Pi$, but the group $\widetilde\Pi$ itself operates by conjugation on the \unsure{set} of splittings of the extension (3), This is in particular \inserted{so} for $\rho$, and we have
\[
\rho(\rho'_g)\overset{\unsure{\mathrm{dfn}}}{=}
\rho\rho'_g\rho^{-1}=\rho(g\rho)\rho^{-1}
=(\rho g\rho^{-1}).\rho=\unsure{g'_{\rho(g)}|\rho'_{\rho(g)}}, \tag{16}
\]
i.e. the operation of $\rho$ on the \unsure{set} of splittings in question, is identified with the operation \deleted{of} induced \deleted{by}\unsure{by|$\rho$} the permutation $\rho$ of $\Pi$ on the \unsure{set} of solutions of (8) $(g\in\Pi)$. Let us note moreover that $g$ and $\rho(g)$ \deleted{\illegible{1}} are $\Pi$-conjugate, for
\[
\rho(g)=\deleted{g}\inserted{\alpha}g\rho(\alpha)^{-1}\qquad\text{with }\alpha=g^{-1}. \tag{17}
\]

Let us note that \deleted{\illegible{4}}\inserted{the operation of $\widetilde\Pi$}\deleted{\unsure{$\rho$}}, on the \unsure{set} of splittings of (3), or again on the \unsure{set} of solutions of (8), is \unsure{always|thus} \unsure{given} by
\[
\Theta_\alpha:g\longmapsto\alpha g\rho(\alpha)^{-1} \tag{10}
\]
where now we assume $\alpha\in\widetilde\Pi$ (and not necessarily $\alpha\in\Pi$). The $\widetilde\Pi$-conjugacy classes are \unsure{reduced|identical} to the $\Pi$-conjugacy classes, \inserted{(as it follows} \unsure{follows} from (17)) \deleted{\unsure{has just seen it.}} In particular, if $\alpha\in\widetilde\Pi$, the element
\[
\alpha\rho(\alpha)^{-1}\qquad(\alpha\in\widetilde\Pi) \tag{18}
\]
of $\widetilde\Pi$ is in $\Pi$, and it \deleted{\illegible{2}} satisfies (8), and